\documentclass[11pt,a4paper]{article}
\usepackage[T1]{fontenc}
\usepackage{lmodern}
\usepackage[margin=29mm]{geometry}
\usepackage{amsmath,amssymb,amsthm}
\usepackage{mathtools}
\usepackage{microtype}
\usepackage{enumitem}
\usepackage{xcolor}
\usepackage{hyperref}
\usepackage{bookmark}
\definecolor{linkblue}{RGB}{25,65,115}
\hypersetup{colorlinks=true,linkcolor=linkblue,citecolor=linkblue,
  urlcolor=linkblue,
  pdftitle={Distributed List Coloring from Lists of Size d(d+1)},
  pdfauthor={AI-generated research note}}
\urlstyle{same}
\setlist{itemsep=3pt,topsep=5pt}
\setlength{\emergencystretch}{2em}
\numberwithin{equation}{section}
\newtheorem{theorem}{Theorem}[section]
\newtheorem{proposition}[theorem]{Proposition}
\newtheorem{lemma}[theorem]{Lemma}
\newcommand{\Nout}{N^{+}}
\newcommand{\Lin}{L_{\mathrm{in}}}
\newcommand{\full}{\mathcal F}
\newcommand{\exc}{\mathcal E}
\newcommand{\normal}{\mathcal N}
\newcommand{\first}{\operatorname{first}}
\newcommand{\outdeg}{\deg^{+}}
\newcommand{\logstar}{\log^{*}}
\newcommand{\LOCAL}{\mathsf{LOCAL}}
\title{Distributed List Coloring\\from Lists of Size
  \texorpdfstring{$d(d+1)$}{d(d+1)}}
\author{AI-generated research note}
\date{30 September 2026}

\begin{document}
\maketitle

\begin{abstract}
Let a finite simple graph be supplied with an orientation of maximum
outdegree at most \(d\geq 2\). We give a uniform deterministic distributed
algorithm that properly colors the graph from arbitrary vertex lists of
size at least \(d(d+1)\) in \(O(\logstar M+d^2)\) rounds, where distinct
vertex identifiers lie in a commonly known namespace of size \(M\).
The algorithm uses only information obtained by reading outgoing
neighbors. Indegree is unrestricted, and the round bound has no
dependence on the universe of list colors. The construction removes one
color from the sufficient list size \(d(d+1)+1\) obtained by an earlier
two-sweep method. Its additional ingredient is a canonical-prefix
argument: the only possible equality obstruction lies in a fixed
independent set, and each vertex of that set has two list colors absent
from the entire input lists of all but one outgoing neighbor.
We give all combinatorial, preprocessing, and round-complexity proofs.
\end{abstract}

\section{Context and contribution}
\label{sec:context}

List coloring assigns each vertex \(v\) a color from its own prescribed
list, with different colors at adjacent vertices. The bound in this
note is a bound on the size of \emph{each list}, not on the total number
of distinct colors appearing anywhere in the input or output.
In particular, different vertices may have disjoint lists.

The starting point is the two-sweep approach to distributed coloring.
Barenboim and Elkin introduced increasing and decreasing color-class
sweeps for defective coloring~\cite{BE09}. Fuchs and Kuhn extended this
approach to list defective coloring~\cite{FK24,FK26}: a vertex first
chooses a set of candidate colors, and a reverse sweep selects its final
color. We use this candidate-set architecture and the restriction to
outgoing neighbors from their oriented setting. The preliminary
reduction of identifier colors to a quadratic palette uses Linial's
separating-set method~\cite{Linial92}, including a finite-field
polynomial construction of Erd\H{o}s, Frankl, and F\"uredi~\cite{EFF85}.
A complete construction and analysis appear in
Section~\ref{sec:preprocessing}.

For a precise comparison, the case with all allowed defects zero
and auxiliary parameter \(\varepsilon=0\) in
Fuchs and Kuhn's theorem~\cite[Theorem~1.1]{FK24}
(also~\cite[Theorem~2.1]{FK26}) permits lists of a common size \(s\)
when
\[
 s>\max\{p,s/p\}\,d
\]
for an integer parameter \(p\geq1\). Taking \(p=d+1\) and
\(s=d(d+1)+1\) satisfies this condition, since
\(\max\{p,s/p\}=d+1\). Their earlier 2024 work therefore already
provides this comparison bound.

We show that \(s=d(d+1)\) suffices for \(d\geq2\).
The proof identifies the exact equality case in a candidate-set
counting argument, reserves two colors at every vertex in a
statically defined exceptional set, and runs the two sweeps on the
remaining vertices. Crucially, the reserved colors are isolated
from the other neighbors' \emph{full input lists}; isolation from
their initial candidate sets would not suffice.
This is a sufficient bound, with no claim of optimality.

\section{Model and main theorem}
\label{sec:model}

The empty graph is trivially colored, so assume \(n=|V|\geq1\).
The graph \(G=(V,E_G)\) is finite and simple, with exactly one supplied
direction on each edge. Write \(v\to u\) for a directed edge and set
\[
 \Nout(v)=\{u:v\to u\},\qquad
 \outdeg(v)=|\Nout(v)|.
\]
A common integer \(d\geq2\) bounds every outdegree. There is no
restriction on indegree, connectedness, or directed cycles. The
outgoing edges at a vertex have distinct local port labels, ordered
locally; no consistency between port orders at different vertices is
required.

Each vertex \(v\) has a distinct identifier in
\([M]=\{1,\ldots,M\}\), where \(M\geq1\) is a common known bound.
It also receives a finite set \(\Lin(v)\) of distinct color names.
Names have finite representations and a common total order, with
consistent equality and comparison available as local computation.
For example, names may be arbitrary nonnegative integers. There is
no supplied bound on their magnitude, on the union of the lists, or on
the lengths of their representations.

We use the synchronous \(\LOCAL\) convention of unrestricted local
computation and message size; see Linial~\cite{Linial92} for the
distributed locality framework. More specifically, the algorithm
obeys the following \emph{forward-information restriction}.
Initially, a vertex knows its own identifier, input list, outgoing
ports, and the common parameters. In one round it can read the
previous-round state of every outgoing neighbor and then update its
own state. Its transition reads no incoming-neighbor state or incoming
incidence information. Thus \emph{forward} describes which vertices may
be observed: for \(v\to u\), \(v\) reads information from \(u\).
It is not a requirement that messages travel from tail to head.
Ordinary two-way communication on the underlying undirected graph
can implement these reads.

All finite computations between rounds are free in this model.
In particular, finite searches used below are not claimed to be
computationally efficient. Uniformity means that one finite program,
given \(d\) and \(M\), performs the stated constructions and terminates
on every legal input; no family of graph-dependent advice is supplied.
All vertices use common phase lengths determined by these parameters.

All logarithms without a subscript are base two; \(\ln\) denotes the
natural logarithm. For \(x\geq1\), \(\logstar x\) is the least number
of iterations of the base-two logarithm needed to reach a value
at most \(1\).

\begin{theorem}
\label{thm:main}
Under the preceding conventions, suppose that
\[
 |\Lin(v)|\geq d(d+1)\qquad\text{for every }v\in V.
\]
There is a uniform deterministic forward-information algorithm that
computes a coloring \(\varphi\) satisfying
\[
 \varphi(v)\in\Lin(v)\quad(v\in V),\qquad
 \varphi(v)\neq\varphi(u)\quad(v\to u)
\]
in \(O(\logstar M+d^2)\) rounds, with absolute constants.
The round bound is independent of the color-name universe.
If \(n=|V|\) and \(M\leq\max\{2,n\}^{C}\) for a fixed constant \(C\),
the bound is \(O(\logstar n+d^2)\).
\end{theorem}

The restriction \(d\geq2\) is necessary for this formula: a cyclically
oriented triangle has outdegree one and has no proper coloring from
the common two-element list.

We first prove the list-coloring stage from a proper ordered coloring,
then construct that ordered coloring from the identifiers.
A \emph{proper ordered coloring} is a map
\(\lambda:V\to[K]\) such that
\(\lambda(v)\neq\lambda(u)\) on every edge.
The order on \([K]\) schedules vertices; it imposes no monotonicity on
the supplied edge directions.

\begin{proposition}
\label{prop:liststage}
Suppose a proper ordered coloring \(\lambda:V\to[K]\) and the common
bound \(K\geq1\) are already available. Under the list-size and
outdegree assumptions of Theorem~\ref{thm:main}, a deterministic
forward-information algorithm computes a proper list coloring in
at most \(2K+4\) additional rounds.
\end{proposition}

\section{The list-coloring stage}
\label{sec:liststage}

\subsection{Canonical prefixes and exceptional vertices}

Fix \(d\geq2\) and put \(k=d(d+1)\). For a finite ordered set \(S\)
and an integer \(r\) with \(0\leq r\leq|S|\), let
\(\first_r(S)\) denote its \(r\) least elements. Retain
\[
 L(v)=\first_k(\Lin(v)),
 \qquad |L(v)|=k.
\]
Every set and count in the setup below is defined from these lists,
the orientation, and \(\lambda\), before either sweep starts.

Let
\[
 \full=\{f:|\Nout(f)|=d
       \text{ and }\lambda(u)>\lambda(f)
       \text{ for every }u\in\Nout(f)\}.
\]
For \(f\in\full\), define its canonical prefix
\begin{equation}
 B(f)=\first_{d+1}(L(f))
     =\first_{d+1}(\Lin(f)).
 \label{eq:prefix}
\end{equation}
The second equality holds because truncation retained the least
\(k\geq d+1\) elements.

A vertex \(e\) is \emph{exceptional} if it has exactly \(d\) outgoing
neighbors, all of them have smaller rank than \(e\), all of them
belong to \(\full\), and
\begin{equation}
 L(e)=\bigsqcup_{f\in\Nout(e)} B(f).
 \label{eq:partition}
\end{equation}
Here the disjoint-union notation requires the \(d\) prefixes to be
pairwise disjoint. Write \(\exc\) for the exceptional set and
\(\normal=V\setminus\exc\) for the set of normal vertices.

\begin{lemma}
\label{lem:independent}
The sets \(\exc\) and \(\full\) are disjoint, every outgoing
neighbor of an exceptional vertex is normal, and \(\exc\) is an
independent set in the underlying undirected graph.
\end{lemma}
\begin{proof}
A vertex of \(\full\) has \(d\) outgoing neighbors of larger rank,
whereas a vertex of \(\exc\) has \(d\) outgoing neighbors of smaller
rank. Since \(d>0\), it cannot have both properties. Every outgoing
neighbor of a vertex of \(\exc\) belongs to \(\full\), hence does
not belong to \(\exc\). Therefore no directed edge has both endpoints
in \(\exc\); every underlying edge has a direction, so the set is
independent.
\end{proof}

A vertex in \(\full\) may have outgoing neighbors in \(\exc\).
In that event, deletions made below may change its eventual candidate
set. The proof will identify that candidate set with \(B(f)\) only
under an additional equality condition. Membership in \(\exc\)
does not assert that a later sweep would actually fail there.

\subsection{Two reserved colors insulated from full input lists}

\begin{lemma}[Prefix isolation]
\label{lem:isolation}
For every \(e\in\exc\), one can deterministically choose a
distinguished outgoing neighbor \(f(e)\) and a two-element set
\(P(e)\) such that
\begin{equation}
 \begin{split}
 P(e)&\subseteq L(e)\cap B(f(e)),\\
 P(e)\cap\Lin(g)&=\varnothing
 \quad\text{for every }g\in\Nout(e)\setminus\{f(e)\}.
 \end{split}
 \label{eq:isolation}
\end{equation}
The choice uses only \(e\)'s list, its outgoing prefixes, and its
local port order.
\end{lemma}
\begin{proof}
Let \(Q(e)=\first_{d+1}(L(e))\).
The \(d\) disjoint prefixes in~\eqref{eq:partition} partition \(L(e)\).
Since \(Q(e)\) has \(d+1\) elements, at least one of these prefixes
contains at least two elements of \(Q(e)\). Among these eligible
neighbors, let \(f(e)\) be the one on the least local outgoing port.
Let \(P(e)\) consist of the two least elements of
\(Q(e)\cap B(f(e))\). This proves the first inclusion.

Fix any other outgoing neighbor \(g\), and put \(q=\max Q(e)\).
Since \(q\) is the \((d+1)\)st least element of \(L(e)\), every
\((d+1)\)-element subset of \(L(e)\) has maximum at least \(q\).
In particular,
\[
 q\leq\max B(g).
\]
For \(x\in P(e)\), we have \(x\leq q\), and disjointness of the
prefixes gives \(x\notin B(g)\).
On the other hand, \eqref{eq:prefix} implies that every element of
\(\Lin(g)\setminus B(g)\) is strictly larger than \(\max B(g)\).
Thus \(x\) belongs neither to \(B(g)\) nor to its complement within
\(\Lin(g)\), proving \(x\notin\Lin(g)\).
\end{proof}

Consequently, once the outgoing neighbors of \(e\) have chosen
arbitrary colors from their own input lists, all but \(f(e)\) avoid
both colors in \(P(e)\). The single color of \(f(e)\) can rule out
at most one of them. This guarantee survives any changes to those
neighbors' candidate sets. Computing the pair requires reading
only the canonical prefixes; the full-list exclusion is a proved
property of the unread list tails.

\subsection{Deleting reserved pairs and choosing candidate sets}

For a normal vertex \(v\), partition its outgoing neighbors into
exceptional neighbors, lower-ranked normal neighbors, and
higher-ranked normal neighbors. Let their respective numbers be
\begin{equation}
\begin{split}
 t(v)&=|\Nout(v)\cap\exc|,\\
 a(v)&=|\{u\in\Nout(v)\cap\normal:\lambda(u)<\lambda(v)\}|,\\
 b(v)&=|\{u\in\Nout(v)\cap\normal:\lambda(u)>\lambda(v)\}|.
\end{split}
\label{eq:counts}
\end{equation}
Properness of \(\lambda\) implies that these are all the possibilities,
and hence
\begin{equation}
 a(v)+b(v)+t(v)=\outdeg(v)\leq d.
 \label{eq:budget}
\end{equation}
Delete the reserved pairs of outgoing exceptional neighbors:
\begin{equation}
 L'(v)=L(v)\setminus
           \bigcup_{e\in\Nout(v)\cap\exc}P(e).
 \label{eq:deleted}
\end{equation}
Overlaps among the deleted pairs can only help, so
\begin{equation}
 |L'(v)|\geq d(d+1)-2t(v).
 \label{eq:deletionbound}
\end{equation}

The increasing sweep processes rank classes \(1,2,\ldots,K\).
Vertices in the same class act simultaneously. A normal vertex \(v\)
chooses the candidate set, or \emph{advertisement},
\begin{equation}
 A(v)=\first_{b(v)+1}
 \left(
 L'(v)\setminus
 \bigcup_{\substack{u\in\Nout(v)\cap\normal\\
                   \lambda(u)<\lambda(v)}} A(u)
 \right).
 \label{eq:advertisement}
\end{equation}
Only advertisements already constructed at smaller ranks occur on
the right. Exceptional vertices do nothing during this sweep.

\begin{lemma}
\label{lem:availability}
Every choice in~\eqref{eq:advertisement} is possible.
In particular, \(|A(v)|=b(v)+1\leq d+1\) for every normal vertex.
\end{lemma}
\begin{proof}
Induct on the processed rank classes. Suppose all required
lower-ranked advertisements have been constructed with the stated
sizes. Fix a vertex \(v\) in the current class, and abbreviate
\(a=a(v)\), \(b=b(v)\), and \(t=t(v)\).
By~\eqref{eq:deletionbound} and the induction hypothesis, the
number of available elements in~\eqref{eq:advertisement} is at least
\begin{equation}
 d(d+1)-2t-a(d+1).
 \label{eq:availablebound}
\end{equation}

First suppose \(t\geq1\). Using \(a\leq d-b-t\), the difference
between~\eqref{eq:availablebound} and the required \(b+1\) is at least
\[
\begin{split}
 d(d+1)-2t-(d-b-t)(d+1)-(b+1)
   &=db+(d-1)t-1\\
   &\geq0,
\end{split}
\]
where the last inequality uses \(d\geq2\).
This includes the tight case \(d=2\), \(t=1\), \(a=1\), \(b=0\).

Next suppose \(t=0\) and \(a\leq d-1\). Put \(s=d-a\geq1\).
The lower bound~\eqref{eq:availablebound} is \(s(d+1)\).
Since \(b\leq s\) and \(s(d+1)\geq s+1\), there are at least \(b+1\)
available elements.

The only remaining case is \(t=0\) and \(a=d\).
Equation~\eqref{eq:budget} forces \(b=0\), so only one element is
needed. Suppose for a contradiction that none is available.
Since \(t=0\), we have \(L'(v)=L(v)\). The \(d\) lower outgoing
advertisements then cover \(L(v)\), giving
\begin{equation}
 d(d+1)=|L(v)|
 \leq
 \left|\bigcup_{u\in\Nout(v)}A(u)\right|
 \leq
 \sum_{u\in\Nout(v)}|A(u)|
 \leq d(d+1).
 \label{eq:saturation}
\end{equation}
Every inequality is an equality. The union contains \(L(v)\)
and has the same cardinality, so it equals \(L(v)\).
The sum of advertisement sizes is \(d(d+1)\), so every one of
the \(d\) advertisements has size \(d+1\).
Equality between union size and the sum of sizes
says that they are pairwise disjoint. Thus they partition \(L(v)\).

For each such neighbor \(u\), the already proved identity
\(|A(u)|=b(u)+1=d+1\) implies \(b(u)=d\).
Applying~\eqref{eq:budget} at \(u\) gives \(a(u)=t(u)=0\).
Consequently \(u\) has exactly \(d\) outgoing neighbors, all normal
and of larger rank; in particular \(u\in\full\).
No pair was deleted from \(L(u)\), and no smaller-ranked
advertisement was excluded at \(u\).
The stipulated least-element rule in~\eqref{eq:advertisement}
therefore gives
\[
 A(u)=\first_{d+1}(L(u))=B(u).
\]
We have proved that \(v\) has \(d\) smaller-ranked outgoing
neighbors in \(\full\) whose canonical prefixes partition \(L(v)\).
These are precisely the conditions defining \(\exc\).
This contradicts \(v\in\normal\).
The putative failure is impossible, and the induction is complete.
\end{proof}

The use of \emph{higher normal} neighbors in \(b(v)\), and the
canonical least-element choice in the increasing sweep, are both
needed in the equality argument. There is no assertion that an
arbitrary vertex of \(\full\) keeps its canonical prefix.

\subsection{The decreasing sweep and exceptional completion}

Process normal vertices in rank order \(K,K-1,\ldots,1\).
At \(v\), choose the least color in
\begin{equation}
 A(v)\setminus
 \{\varphi(u):u\in\Nout(v)\cap\normal,\
                    \lambda(u)>\lambda(v)\}.
 \label{eq:normalcolor}
\end{equation}
There are \(b(v)\) higher-ranked normal outgoing neighbors and
\(|A(v)|=b(v)+1\), so the choice is possible.

After this sweep, all normal vertices have final colors.
For each exceptional vertex \(e\), simultaneously set
\begin{equation}
 \varphi(e)=\min\bigl(P(e)\setminus\{\varphi(f(e))\}\bigr).
 \label{eq:exceptioncolor}
\end{equation}
Lemma~\ref{lem:independent} guarantees that \(f(e)\) is normal,
so its color is already defined. A single color cannot exhaust the
two-element set \(P(e)\).

\begin{lemma}
\label{lem:correctness}
The resulting \(\varphi\) is a proper coloring and
\(\varphi(v)\in\Lin(v)\) at every vertex.
\end{lemma}
\begin{proof}
For a normal vertex, the construction gives
\[
 \varphi(v)\in A(v)\subseteq L'(v)\subseteq L(v)\subseteq\Lin(v).
\]
For an exceptional vertex, it gives
\(\varphi(e)\in P(e)\subseteq L(e)\subseteq\Lin(e)\).
It remains to verify every directed edge \(v\to u\).

If \(v,u\in\normal\) and \(\lambda(u)<\lambda(v)\), the increasing
sweep chose \(A(v)\) disjoint from \(A(u)\), so their final colors
differ. If \(v,u\in\normal\) and \(\lambda(u)>\lambda(v)\),
the decreasing choice at \(v\) explicitly excluded \(\varphi(u)\).
There is no equal-rank case.

If \(v\in\normal\) and \(u\in\exc\), then
\(A(v)\subseteq L'(v)\) is disjoint from \(P(u)\) by~\eqref{eq:deleted}.
Since \(\varphi(u)\in P(u)\), the edge is proper.

If \(v\in\exc\) and \(u\in\normal\), the choice
in~\eqref{eq:exceptioncolor} excludes \(\varphi(u)\) when \(u=f(v)\).
Otherwise Lemma~\ref{lem:isolation} gives
\(P(v)\cap\Lin(u)=\varnothing\), whereas
\(\varphi(u)\in\Lin(u)\) and \(\varphi(v)\in P(v)\).
Again the colors differ.

There are no edges with both endpoints exceptional, by
Lemma~\ref{lem:independent}. These cases cover every underlying edge.
\end{proof}

\subsection{Communication schedule}
\label{sec:schedule}

Here is a complete schedule under the read-previous-state convention
of Section~\ref{sec:model}. Vertices retain any state needed in later
rounds.
\begin{enumerate}
 \item In the first setup round, each vertex reads its outgoing
 neighbors' ranks and determines membership in \(\full\).
 If it belongs to \(\full\), it computes its canonical prefix locally.
 \item In the second round, it reads its outgoing neighbors'
 membership in \(\full\), together with the prefixes of those in
 \(\full\). It can now test
 membership in \(\exc\); if exceptional, it chooses \(f(e)\)
 and \(P(e)\).
 \item In the third round, it reads its outgoing neighbors'
 exceptional flags, together with the pairs of those in \(\exc\).
 Every normal vertex can
 now determine \(a(v),b(v),t(v)\) and \(L'(v)\).
 \item Allocate \(K\) rounds to the increasing sweep, one per
 rank class. At its scheduled round a vertex reads the stored
 advertisements of its smaller-ranked outgoing normal neighbors.
 \item Allocate \(K\) rounds to the decreasing sweep.
 At its scheduled round a vertex reads the final colors of its
 higher-ranked outgoing normal neighbors.
 \item One final round gives each exceptional vertex its
 distinguished outgoing neighbor's final color, allowing
 \eqref{eq:exceptioncolor}.
\end{enumerate}
The total is \(3+K+K+1=2K+4\) rounds.
Adjacent vertices have different ranks, and every dependency in
each sweep points to a rank class already processed in that sweep.
This remains true on directed cycles. No distinguished neighbor
needs to learn which incoming vertices selected it.
All communicated objects are finite and every operation is
deterministic, using ordered lists or local port order to break ties.
Together with Lemmas~\ref{lem:availability} and~\ref{lem:correctness},
this proves Proposition~\ref{prop:liststage}.

\section{Self-contained construction of the ordered coloring}
\label{sec:preprocessing}

We now supply the preprocessing needed for the main theorem.
The separating-set transformation is Linial's color-reduction
method~\cite[Section~4]{Linial92}; the polynomial sets used in its
last step come from Erd\H{o}s, Frankl, and F\"uredi~\cite[Example~3.2]{EFF85}.
We prove the required properties here, including a round bound
uniform in \(d\), and specify finite deterministic searches so that
no nonuniform advice is required. These computations use identifiers
and temporary labels, not the input list colors.

\begin{proposition}
\label{prop:preprocessing}
From identifiers in the common namespace \([M]\), one can compute
a proper ordered coloring \(\lambda:V\to[K]\) with
\[
 K<256(d+1)^2
\]
in at most \(2\logstar M+3\) forward-information rounds.
\end{proposition}

\subsection{A separating family and one-round reduction}

Put \(u=d+1\) and \(A=8u^2\). For any integer \(h\geq2\), let
\[
 m=A\lceil\log h\rceil.
\]
We claim that there are sets \(S_1,\ldots,S_h\subseteq[m]\) such that
\begin{equation}
 S_i\setminus\bigcup_{j\in J}S_j\neq\varnothing
 \quad\text{whenever }i\in[h],\
 J\subseteq[h]\setminus\{i\},\ |J|\leq d.
 \label{eq:separating}
\end{equation}

To prove existence, independently include each point in each set
with probability \(1/u\). For a fixed pair \((i,J)\), a given point
belongs to \(S_i\) and to none of the sets indexed by \(J\) with
probability at least
\[
 \frac1u\left(1-\frac1u\right)^d
 =\frac1u\left(1+\frac1d\right)^{-d}
 \geq\frac1{eu}.
\]
The last inequality follows from \(\ln(1+x)\leq x\) for \(x\geq0\).
Independence across points and \(1-z\leq e^{-z}\) imply that
the probability that this pair has no witness is at most
\begin{equation}
 \exp\!\left(-\frac{m}{eu}\right)
 \leq \exp(-2u\ln h)=h^{-2u}.
 \label{eq:failureprobability}
\end{equation}
Indeed, \(m/(eu)=(8/e)u\lceil\log h\rceil\geq2u\ln h\),
using \(e<3\) and \(\log h\geq\ln h\).

There are at most \(h^{d+1}=h^u\) pairs \((i,J)\).
For fixed \(i\), encode \(J\) by its increasing list of elements,
followed by copies of \(i\) until the sequence has length \(d\).
This is an injection into \([h]^d\), since \(i\notin J\).
The union bound and~\eqref{eq:failureprobability} put the probability
of any failed pair at most \(h^{-u}<1\). A family
satisfying~\eqref{eq:separating} therefore exists.

To select a common family deterministically, order all \(h\)-tuples
of subsets of \([m]\) lexicographically by their incidence bit
strings. Enumerate them until all the finitely many conditions
in~\eqref{eq:separating} hold. The existence proof guarantees
termination, and the common parameters guarantee that every vertex
finds the same family.

Suppose the current proper coloring has labels in \([h]\).
A vertex of label \(i\) reads its outgoing neighbors' labels in
one round and chooses the least element of
\begin{equation}
 S_i\setminus
 \bigcup_{w\in\Nout(v)}S_{\mathrm{old}(w)}.
 \label{eq:onestep}
\end{equation}
Neighbor labels differ from \(i\); repeated labels at different
neighbors cause no problem because there are at most \(d\) distinct
ones. The choice is possible by~\eqref{eq:separating}.
For an edge \(v\to w\), the new label of \(v\) is outside
\(S_{\mathrm{old}(w)}\), whereas the new label of \(w\) belongs to
that set. Hence the new coloring is proper and has palette \([m]\).

\subsection{Bounding the number of reductions uniformly}

Let
\begin{equation}
 H_0=16A\lceil\log(4A)\rceil,\qquad
 f(x)=A\lceil\log x\rceil\quad(x\geq2).
 \label{eq:threshold}
\end{equation}
Starting with the identifier bound \(h=M\), apply the one-round
reduction while \(h>H_0\), replacing the common bound by \(f(h)\).
If \(M=1\) or already \(M\leq H_0\), no such round is necessary.
In the former case, distinct identifiers imply that there is at most
one vertex. In every case the current proper labels can be viewed
as elements of \([H_0]\) once the loop has ended.

For completeness, we prove that the loop takes at most
\(2\logstar M+2\) rounds, with constants independent of \(d\).
Set
\[
 T=(4A)^4.
\]
For \(x\geq2^T\), write \(y=\log x\geq T\).
As \(\lceil y\rceil\leq2y\) for \(y\geq1\), monotonicity gives
\begin{equation}
 f(f(x))
 \leq A\bigl(\log(2Ay)+1\bigr)
 \leq y=\log x.
 \label{eq:tworeductions}
\end{equation}
Here is an explicit verification of the second inequality.
For \(y\geq T\), the function
\[
 g(y)=y-A\bigl(\log(2Ay)+1\bigr)
\]
has derivative \(g'(y)=1-A/(y\ln2)>0\).
Moreover,
\[
\begin{split}
 A\bigl(\log(2AT)+1\bigr)
 &=A\bigl(\log(2A)+4\log(4A)+1\bigr)\\
 &\leq18A^2+A
 <256A^4=T,
\end{split}
\]
where \(\log z\leq z\) for \(z\geq1\), and \(A\geq32\).
Thus \(g(T)>0\), proving~\eqref{eq:tworeductions}.

As long as a pair of rounds starts with bound at least \(2^T\),
either the loop stops during that pair, or its new bound is at
most the logarithm of the bound at the start of the pair.
More explicitly, let \(x_j\) be the bound after \(2j\) executed
rounds. If the first \(j+1\) pairs are executed and each starts
with bound at least \(2^T\), then \(x_{j+1}\leq\log x_j\);
induction gives
\(x_j\leq\log^{(j)}M\), where \(\log^{(j)}\) denotes \(j\)
successive logarithms. Thus after at most \(\logstar M\) pairs
the loop has stopped or its bound has dropped below \(2^T\);
otherwise its bound would eventually be at most \(1\).

Once the bound is below \(2^T\), at most two more rounds are needed.
Indeed, by monotonicity the first new bound is at most
\(f(2^T)=AT\), and the second is at most
\[
\begin{split}
 f(AT)
 &\leq A\bigl(\log(AT)+1\bigr)\\
 &\leq 6A\log(4A)
 <16A\lceil\log(4A)\rceil=H_0.
\end{split}
\]
The middle inequality follows from
\(\log(AT)=\log A+4\log(4A)\leq5\log(4A)\) and
\(\log(4A)\geq1\). The loop stops earlier if it reaches \(H_0\)
at an intermediate round. This proves the claimed bound.

\subsection{A final polynomial reduction to a quadratic palette}

Choose \(q\) to be the smallest power of two at least \(8u\), so
\begin{equation}
 8u\leq q<16u.
 \label{eq:q}
\end{equation}
We use the polynomial graph construction of
Erd\H{o}s, Frankl, and F\"uredi~\cite[Example~3.2]{EFF85},
as in Linial's color reduction~\cite[Lemma~4.2]{Linial92}.

We will need a field \(\mathbb F\) of order \(q\).
We recall its existence explicitly. Write \(q=2^r\).
Over the two-element field, form a splitting field of \(X^q-X\).
Such a field can be constructed by repeatedly adjoining a root of
an irreducible factor of degree at least two until all factors
are linear: the quotient of a polynomial ring over a
field by an irreducible polynomial is a field, because the Euclidean
algorithm supplies an inverse modulo that polynomial for every
nonzero residue class. At each step a further factor becomes linear,
so a finite succession of such extensions suffices.

The derivative of \(X^q-X\) in characteristic two is \(-1\), so
in its splitting field the polynomial has exactly \(q\) distinct
roots. These roots form a field: they contain \(0\) and \(1\), and
are closed under addition and multiplication because repeated
squaring gives \((a+b)^q=a^q+b^q\) and
\((ab)^q=a^qb^q\). Additive inverses cause no problem in
characteristic two. For a nonzero root \(a\), the identity
\((a^{-1})^q=(a^q)^{-1}=a^{-1}\) gives closure under inverses.
Thus their set is a field with exactly \(q\) elements.

For a uniform common representation, enumerate pairs of binary
operation tables on the ordered set \([q]\), with prescribed
distinct labels for \(0\) and \(1\), and choose the first pair
satisfying the field axioms. This finite search terminates by the
existence argument. Field elements are ordered by their labels;
this is merely an order for deterministic choices, not an order
compatible with field arithmetic.

There are \(q^3\) polynomials of degree at most two over this field,
indexed by their triples of coefficients. This is enough to
assign a distinct polynomial \(P_i\) to every temporary color
\(i\in[H_0]\). To check this, for integer \(u\geq2\) the inequality
\(u\leq2^{u-1}\) follows by induction, and hence
\[
 \left\lceil\log(32u^2)\right\rceil
 \leq 2u+3\leq4u.
\]
Since \(A=8u^2\), we obtain
\begin{equation}
 H_0
 =128u^2\left\lceil\log(32u^2)\right\rceil
 \leq512u^3
 \leq q^3.
 \label{eq:polynomialcount}
\end{equation}
Use coefficient lexicographic order to make the assignment common
and deterministic.

For each \(i\in[H_0]\), put
\[
 H_i=\{(a,P_i(a)):a\in\mathbb F\}\subseteq\mathbb F^2.
\]
Each \(H_i\) has \(q\) elements. Distinct \(H_i,H_j\) intersect in
at most two points, because \(P_i-P_j\) is a nonzero polynomial of
degree at most two. The root bound used here follows by division
by \(X-a\): a root \(a\) gives a factor \(X-a\), and induction on
degree bounds the number of distinct roots by the degree.

At most \(d\) outgoing neighbors can therefore cover at most
\(2d<q\) points of a vertex's own set \(H_i\).
In one more round, a vertex of temporary color \(i\) chooses the
least pair in its \(H_i\) outside the sets corresponding to its
outgoing neighbors. This is precisely the transformation
in~\eqref{eq:onestep} with the \(H_i\) in place of the \(S_i\),
so its correctness proof applies verbatim.
Order \(\mathbb F^2\) lexicographically by the common field labels
and relabel its elements by \([q^2]\).
The result is a proper ordered coloring with
\[
 K=q^2<256u^2=256(d+1)^2.
\]
The previous loop took at most \(2\logstar M+2\) rounds, and this
last reduction takes one. This proves
Proposition~\ref{prop:preprocessing}.

\section{Proof of the main theorem}

Run Proposition~\ref{prop:preprocessing}, then apply the algorithm
of Proposition~\ref{prop:liststage} using its output ranks.
The list stage preserves membership in the actual input lists
and separates every edge, by Lemma~\ref{lem:correctness}.
Its composition with preprocessing takes at most
\[
 (2\logstar M+3)+(2K+4)
 <2\logstar M+512(d+1)^2+7
 =O(\logstar M+d^2)
\]
rounds. Under a polynomial identifier bound with fixed exponent,
\(\logstar M=O(\logstar n+1)\), giving the stated bound in \(n\).

All searches in preprocessing have finite domains, explicitly
testable conditions, and proved existence of a successful object.
All other choices are least elements of proved nonempty finite
sets, with local port order resolving the distinguished-neighbor
choice. The common parameters determine every phase length.
These facts give one terminating deterministic program for all
values of \(d\) and \(M\).

Every communication step reads only outgoing-neighbor states.
The algorithm uses no incoming-degree bound and no acyclicity
assumption. In the list stage it compares and transmits only
finitely many input names and their subsets; it never enumerates
a global color universe. Longer representations can increase
message size and local computation, which the stated model does
not charge as rounds. This completes the proof of
Theorem~\ref{thm:main}.
\qed

\begin{thebibliography}{9}

\bibitem{BE09}
Leonid Barenboim and Michael Elkin.
\newblock Distributed \((\Delta+1)\)-coloring in linear (in \(\Delta\)) time.
\newblock In \emph{Proceedings of the 41st Annual ACM Symposium on
Theory of Computing (STOC)}, pages 111--120, 2009.
\newblock \href{https://doi.org/10.1145/1536414.1536432}
{doi:10.1145/1536414.1536432}.

\bibitem{EFF85}
Paul Erd\H{o}s, Peter Frankl, and Zolt\'an F\"uredi.
\newblock Families of finite sets in which no set is covered by the
union of \(r\) others.
\newblock \emph{Israel Journal of Mathematics}, 51(1--2):79--89, 1985.
\newblock \href{https://doi.org/10.1007/BF02772959}
{doi:10.1007/BF02772959}.

\bibitem{FK24}
Marc Fuchs and Fabian Kuhn.
\newblock Simpler and more general distributed coloring based on simple
list defective coloring algorithms.
\newblock \href{https://arxiv.org/abs/2405.04648v1}
{arXiv:2405.04648v1}, 2024.
\newblock A brief announcement appeared as
\emph{Brief Announcement: Simpler and More General Distributed Coloring
Based on Simple List Defective Coloring Algorithms},
in \emph{Proceedings of the 43rd ACM Symposium on Principles of Distributed
Computing (PODC)}, pages 425--428, 2024.
\newblock \href{https://doi.org/10.1145/3662158.3662808}
{doi:10.1145/3662158.3662808}.
\newblock The theorem cited in this note is from the full preprint.

\bibitem{FK26}
Marc Fuchs and Fabian Kuhn.
\newblock Greedy-like defective coloring: Distributed algorithms and
applications.
\newblock \href{https://arxiv.org/abs/2608.02386v2}
{arXiv:2608.02386v2}, 10 September 2026.
\newblock \href{https://doi.org/10.48550/arXiv.2608.02386}
{doi:10.48550/arXiv.2608.02386}.

\bibitem{Linial92}
Nathan Linial.
\newblock Locality in distributed graph algorithms.
\newblock \emph{SIAM Journal on Computing}, 21(1):193--201, 1992.
\newblock \href{https://doi.org/10.1137/0221015}
{doi:10.1137/0221015}.

\end{thebibliography}
\end{document}
